\section{Introduction}
\label{sec:introduction}

Salivary gland fluid secretion has served as a particularly useful example of how transporters, ion channels, intracellular signalling and water movement can be assembled into a quantitative epithelial model. In the accepted picture, basolateral transport concentrates intracellular chloride above electrochemical equilibrium. Agonist induced calcium signals then open apical chloride channels and potassium channels, creating a transepithelial salt gradient that drives water into the acinar lumen \citep{Melvin2005,Frizzell2012}. NKCC1 supplies most of the chloride loading and is essential for normal secretion \citep{Evans2000}. A second, bicarbonate dependent route combines sodium/proton exchange with chloride/bicarbonate exchange and has long been thought to support both chloride uptake and intracellular pH regulation \citep{Nauntofte1992,Melvin2005}.

The identity of the relevant chloride/bicarbonate exchanger was unsettled for many years. AE2 is widely expressed and was the obvious candidate. Experiments in mouse submandibular gland instead showed that deletion of AE4 reduced stimulated secretion by approximately one third, whereas deletion of AE2 had little effect \citep{PenaMunzenmayer2015}. AE4 was subsequently shown to be an electroneutral monovalent cation dependent chloride/bicarbonate exchanger \citep{PenaMunzenmayer2016}. These findings posed a simple but nontrivial question: why should AE4 have such strong control over secretion when NKCC1 remains the dominant chloride loader and AE2 is available as another chloride/bicarbonate exchanger?

Our 2018 model was constructed to answer that question \citep{VeraSiguenza2018}. It extended earlier salivary secretion models \citep{Gin2007,Palk2010,Maclaren2012} by including AE2 and a cation coupled AE4 cycle. The model reproduced a negligible AE2 knockout effect and an approximately 24\% reduction in flow after AE4 deletion. It suggested that AE4 was difficult to replace because its chloride uptake was coupled to the export of sodium and potassium. Removing that cation efflux changed sodium, potassium, pump and NHE1 activity in ways that AE2 could not reproduce. The broad biological prediction was attractive: AE4 was not merely another route for chloride entry; its cation dependence gave it a distinct network role.

The model also contained clues that its mechanism was incomplete. First, the simulated knockout phenotype developed within roughly two minutes, much faster than the gradual experimental divergence over approximately ten minutes. Second, the acid--base system was deliberately minimal. Carbon dioxide hydration, NHE1, AE2 and AE4 were included, but no additional proton or bicarbonate transporter was present. This simplification was sufficient for the original calibration, but its structural consequences were not explored. Third, several transporter laws and regulatory assumptions were inherited from earlier models or represented by simple concentration dependent expressions. Later work improved the spatial description of calcium signalling and fluid transport \citep{VeraSiguenza2019Fluid,Pages2019,VeraSiguenza2020,Sneyd2021}, but did not revisit the central AE4 mechanism.

The experimental picture has also changed. AE4 activity in salivary acinar cells is increased by beta adrenergic stimulation through cAMP dependent protein kinase, with serine 173 implicated in the response \citep{PenaMunzenmayer2021}. More recently, molecular and functional experiments showed that sodium and potassium do not interact identically with human AE4 and proposed several possible transport cycles, including sodium and chloride moving inward while potassium and bicarbonate or carbonate move outward \citep{Catalan2025}. The cation dependence anticipated by the original model has therefore received direct molecular support. The exact whole cycle stoichiometry, however, remains unresolved.

These developments allow a more demanding question than the one posed in 2018. Rather than asking whether a model can be made to reproduce the AE4 knockout phenotype, we ask which parts of the original explanation survive a source constrained reconstruction of the whole cell. We separate this into four steps, summarised in \cref{fig:model-evolution}. First, we rebuild wild type transport with explicit amount, charge, carbon, alkalinity and water balances. Second, we determine whether substantial stimulated AE4 chloride loading is compatible with physiological pH. Third, we freeze the resulting model and test whether equal cation routing or the specific stoichiometric classes suggested by recent experiments recover the held out secretion phenotype. Finally, we ask what additional coupling would be mathematically sufficient if stoichiometry alone fails.

This approach changes the role of the phenotype. The approximately 35\% secretion reduction is not used to calibrate the wild type model or to choose among the predeclared AE4 source classes. It is revealed only after those predictions are frozen. A separate inverse construction is then clearly labelled as target selected. The distinction is important. A model can reproduce one macroscopic output for many reasons; the accompanying sodium, chloride, pH, volume, transporter fluxes and time course determine whether the proposed mechanism is credible.

\begin{figure}[t]
    \centering
    \includegraphics[width=\textwidth]{figure1_model_evolution.pdf}
    \caption{The development of the modelling question. (a) The 2018 model included NKCC1, the sodium/potassium pump, NHE1, AE2, AE4, calcium activated channels, carbon dioxide exchange and osmotic water transport, but no sodium/bicarbonate cotransporter. The AE4 knockout phenotype was attributed mainly to altered cation balance. (b) The present reconstruction separates wild type repair, blinded stoichiometric mechanism tests and a target selected constructive comparison.}
    \label{fig:model-evolution}
\end{figure}
